{\Huge\bfseries\color{ink} Molecule Lab}\par
{\Large An academic chemistry laboratory in virtual reality}\par
\textbf{Treatment report / working design draft}\hfill 29 September 2026

\section*{Intent and audience}
A single-user, offline Quest 3 environment for introductory chemistry students. Learners assemble authored molecules, inspect spatial relationships, scale models, twist ethane, split and rejoin simplified radical fragments, and activate four guided, balanced reactions. The purpose is spatial understanding and confident interaction, not general molecular dynamics.

The new environment gives these tasks a recognizable academic setting: a central laboratory table, finite atom stock in a nearby shelf, and subdued scientific equipment around the room. The learner's attention should move naturally from stock, to bonding target, to completed structure.

\section*{Learning and experience goals}
\begin{itemize}
\item Recognize the elements and constituent atom counts of eleven molecules, including benzene, ethene, ethyne, alcohols and carbon dioxide.
\item Relate authored geometry to bond direction; distinguish rotation about a bond from deformation and bond breaking.
\item Observe conservation of atoms in $2\mathrm{H}_2+\mathrm{O}_2\rightarrow2\mathrm{H}_2\mathrm{O}$.
\item Identify manipulable objects without trial-and-error grabbing of scenery.
\item Complete a short sequence without restarting or becoming stuck after an error.
\end{itemize}

\section*{Core loop}
\begin{center}
\begin{tikzpicture}[node distance=5mm]
\node[state,text width=25mm] (a) {Choose recipe};
\node[state,text width=25mm,right=of a] (b) {Take atom};
\node[state,text width=25mm,right=of b] (c) {Align and release};
\node[state,text width=25mm,right=of c] (d) {Inspect / react};
\draw[flow] (a)--(b);\draw[flow] (b)--(c);\draw[flow] (c)--(d);
\draw[flow] (d.south)--++(0,-.5)-| node[pos=.25,below,font=\small]{Recycle and continue} (a.south);
\end{tikzpicture}
\end{center}

\section*{Design basis and boundaries}
The course evaluation slides emphasize clear goals, immediate feedback, scaffolding, control, affordance, faithful viewpoints and consistent departures from realism \cite{gain}. These principles guide the design choices in this report; they are not evidence of achieved usability.

Hands are the primary input with Touch Plus fallback. Artificial locomotion is excluded from this iteration: the required work is arranged around a reachable station. The room supplies context, not an exploration objective. The finite inventory is a visible resource model, not a timed scarcity challenge.

\textbf{Status language.} ``Implemented'' means present in the software; ``verified locally'' means checked in the editor/build; ``validated on headset'' requires recorded device evidence. No usability, learning or performance results are invented here.
